\documentclass[a4paper]{article}

%% Language and font encodings
\usepackage[english]{babel}
\usepackage[utf8x]{inputenc}
\usepackage[T1]{fontenc}
\usepackage{booktabs}

%% Sets page size and margins
\usepackage[a4paper,top=3cm,bottom=2cm,left=3cm,right=3cm,marginparwidth=1.75cm]{geometry}

%% Useful packages
\usepackage{amsmath}
\usepackage{graphicx}
\usepackage[colorinlistoftodos]{todonotes}
\usepackage[colorlinks=true, allcolors=blue]{hyperref}

\title{HLTH301 Assignment at UC}
\author{Write your Student ID here}
\date{}

\begin{document}
\maketitle

\begin{abstract}
Write a short summary of the third assignment here. In the summary include the question you answered, a few sentences to describe the nature of the problem. The problem statement or introduction should introduce the health issue that you are answering and what is known about the health problem. Also indicate why answering the question is important. Then write a few sentences about what databases you searched, and what were you search criteria (you should describe that in words, do not include any tables or figure here). Then write the results you obtained. Here, write briefly how many articles you originally retrieved and how many articles (five, but if not five then how many) articles you finally considered to be good enough to answer the question and which article you found best out of these five and why. Then briefly describe the nature of the article and write a few sentences as to why and how this article answered the main question you set out to answer in this assignment.

Typically the summary will be one paragraph, but you can include more than one paragraph to break up the thoughts and ideas presented in your assignment. Try to be brief and capture the essential point in the assignment you write. Do not include any references or citations in this section. 

\end{abstract}

\section*{Introduction}

the introduction section should briefly introduce the topic or the question that you are answering in this assignment. The introduction section should consist of brief statements about what is the health problem that is being addressed, the intervention that is being addressed, and what is known about the health condition and the intervention or the exposure variable that is being considered. The aim of this section is to introduce the real discussion that will then be built.

This section can contain more than one paragraph, depending on the ideas introduced. For intervention related questions (most questions in the Assignment are related to interventions), you need to describe both the intervention and the disease condition in the introduction section. Also, discuss why answering this question is important or what will be achieved in answering this question. For example, for the second question “How effective is cochlear implant for children with congenital hearing loss?”, you will need to discuss three points, (1) what is meant by congenital hearing loss (what is it, how frequent is it, who suffer from congenital hearing loss, etc), (2) what is cochlear implant,and (3) why correction of hearing loss in children is important and how does cochlear implant may help 

\section*{Search Methods}

In the search methods, you should write clearly what search strategy you used to find the first five articles and eventually narrowed down to the article that you reviewed. Write the search methods in a way that makes it easy for anyone to replicate the searches you conducted. We shall use the same search to test whether we arrive at the same search results. If we do not arrive at the same set of results, then you will lose mark on this aspect. We recommend that you use the following template of the table that you should use to do this section.


\begin{table}[ht]
\centering
\caption{Search Results}
\label{tab:srch}
\begin{tabular}{@{}ll@{}}
\toprule
Search Algorithm           & Number of Articles Retrieved                                      \\ \midrule
Use your search terms here & Write the number of articles you retrieved using this search term \\ \bottomrule
\end{tabular}
\end{table}

In the search history you must include the restrictions you used. For instance, if you searched the last five years of data, then the beginning and end date of the years searched must be included. 

After detailing the search methods you should indicate the rationale for selecting the best five articles and indicate the best article (or one article) that you selected out of the five articles selected. 


\section*{Results}

First, include a summary table like this. You will have to use your summary table from GradePro web app, see Figure \ref{fig:gradepro}. 

\begin{figure}
\centering
\includegraphics[width=0.8\textwidth]{gradepro}
\caption{\label{fig:gradepro}Results from Gradepro}
\end{figure}

\section*{Discussion}

Finally, write the discussion section. The discussion section must contain the following three points: a short summary of the results (only high impact points, perhaps a sentence or two at the most to highlight just what you wrote), (2) how and why this article answered the question — that is, the overall quality of the article, how this article addressed the significant biases and therefore had passed your quality criteria, and what ways does the article addresses the key question. 

Also make sure that your article is well referenced, and references should follow the convention of APA 6th edition (although we are not particularly choosy about APA 6th edition, as long as you clearly attribute your articles to authors and a citation is identifiable that should be OK). 

\subsection*{Guidance: How to add Citations and a References List}

You can upload a \verb|.bib| file containing your BibTeX entries, created with JabRef; or import your \href{https://www.overleaf.com/blog/184}{Mendeley}, CiteULike or Zotero library as a \verb|.bib| file. You can then cite entries from it, like this: \cite{greenwade93}. Just remember to specify a bibliography style, as well as the filename of the \verb|.bib|.

You can find a \href{https://www.overleaf.com/help/97-how-to-include-a-bibliography-using-bibtex}{video tutorial here} to learn more about BibTeX.



\bibliographystyle{alpha}
\bibliography{sample}

\end{document}